%% Computerized Medical Imaging and Graphics -- Elsevier
%% Reference cohort composition in distributional evaluation of medical
%% image generators.
\documentclass[preprint,12pt,authoryear]{elsarticle}

\usepackage{amsmath,amssymb,amsfonts}
\usepackage{graphicx}
\usepackage{booktabs}
\usepackage{multirow}
\usepackage{array}
\usepackage{url}
\usepackage[hidelinks]{hyperref}
\usepackage{xcolor}
\usepackage{enumitem}
\usepackage{subcaption}

\journal{Computerized Medical Imaging and Graphics}

\usepackage{pifont}
\newcommand{\cmark}{\ding{51}}
\newcommand{\xmark}{\ding{55}}
\newcommand{\mmd}{\mathrm{MMD}}
\newcommand{\Lb}[1]{\ensuremath{\mathrm{L}#1}}

\begin{document}

\begin{frontmatter}

\title{Reference cohort composition confounds distributional evaluation of
       medical image generators: a depth-resolved audit of radiology
       foundation-model features}

\author[pdn]{Sathiyamohan Nishankar}
\ead{e17230@eng.pdn.ac.lk}
\address[pdn]{Department of Computer Engineering, University of Peradeniya,
              Peradeniya, Sri Lanka}

\begin{abstract}
Distributional metrics such as the Fr\'echet Inception Distance compare a set of
generated images against a finite set of real ones. In natural-image benchmarks
the real reference is a collection of independent photographs, so its
composition is rarely examined. Medical imaging violates that assumption twice
over: adjacent slices from one patient are near-duplicates, and multi-site
cohorts carry systematic acquisition differences. We show that these two
properties, not generator quality, can dominate the reported number. On BraTS
brain MRI, the maximum mean discrepancy (MMD$^2$) between two \emph{real}
patient cohorts drawn from the same dataset reaches $0.096$, exceeding the
$0.080$ measured between a real cohort and the output of an unconditional DDPM;
a reference built by taking the first $500$ filenames in sorted order --- a
common convenience --- selects a $17$-subject block that is systematically
atypical of the cohort, and inflates the reported score by $1.9\times$ relative
to a subject-diverse reference of identical size. We propose replacing the conventional permutation
test, which asks whether real and generated samples share a distribution and
is therefore rejected in every configuration we tested, with an \emph{empirical inter-cohort reference distribution} built from
distances between disjoint real patient blocks; this reports whether a
generator differs from the reference by more than routine inter-cohort
variation, and it makes the dependence of any reported score on cohort size
explicit. Using this protocol we then audit which encoder depth of a frozen
radiology vision transformer (RadioDINO-s16) supports which evaluation task, at
all twelve blocks. Depth is strongly task-dependent and the final block is not
the best choice for any task we measure: distributional discriminability peaks
at a mid-to-late plateau (\Lb{7}--\Lb{9}) and is $22\%$ lower at \Lb{12};
near-duplicate detection improves monotonically with depth; degradation
monotonicity fails at \Lb{1}--\Lb{2}; and $k$-nearest-neighbour precision/recall
coverage fails its own mode-drop validation at \emph{every} depth, with recall
rising rather than falling as diversity is removed. For per-image novelty
scoring, the choice of scoring rule matters more than depth: a shrinkage
Mahalanobis distance reaches AUC $0.972$ where distance to the real centroid
reaches $0.846$. We release the evaluation code and report a concrete protocol:
report the number of distinct subjects behind every reference set, calibrate
against a cohort null, audit read-out depth rather than defaulting to the final block, and do not
report $k$-NN coverage in this feature space without validating it.
\end{abstract}

\begin{keyword}
generative models \sep medical image synthesis \sep evaluation metrics \sep
maximum mean discrepancy \sep dataset bias \sep foundation models
\end{keyword}

\end{frontmatter}

% =============================================================================
\section{Introduction}
\label{sec:intro}
% =============================================================================

Generative models are increasingly proposed as a response to data scarcity in
medical imaging, where privacy constraints, annotation cost and disease rarity
limit dataset size \citep{kebaili2023deep, kazerouni2023diffusion}. Whether a
synthetic dataset is fit for a downstream purpose depends on being able to
measure how far it departs from real data, and the measurement instruments in
common use were designed for a different setting. The Fr\'echet Inception
Distance (FID) \citep{heusel2017gans} and Kernel Inception Distance (KID)
\citep{binkowski2018demystifying} embed images with an ImageNet classifier and
summarise the gap between two feature distributions as a scalar. Both have been
transplanted into medical image synthesis largely unchanged.

Most methodological criticism of this transplant has focused on the
\emph{embedding}: an ImageNet backbone encodes texture statistics that need not
correspond to radiological structure, motivating replacements built on
domain-pretrained encoders \citep{dohmen2025similarity, woodland2024feature}.
That criticism is well founded, and we take a domain-pretrained encoder as our
starting point. But it leaves a second assumption unexamined. Every
distributional metric is a comparison against a \emph{finite real reference
set}, and its value is a property of that reference as much as of the generator.
In natural-image benchmarks the reference is tens of thousands of independent
photographs, so its composition is close to the population it stands for.
Medical references violate this in two specific ways:

\begin{enumerate}[leftmargin=*, topsep=2pt, itemsep=2pt]
\item \textbf{Slices are not independent samples.} A reference of $N$ axial
      slices may come from $N$ patients or from a dozen; neighbouring slices of
      one volume are near-duplicates, so the effective sample size is closer to
      the number of subjects than to $N$.
\item \textbf{Cohorts are heterogeneous.} Public collections such as BraTS
      aggregate several contributing institutions with different scanners and
      protocols. Which subjects enter the reference therefore carries a site
      signature.
\end{enumerate}

\noindent This paper asks how much of a reported metric value those two
properties account for, and finds that on BraTS they can account for more than
the generator does. Two real patient cohorts drawn from the same dataset can
differ by more, in the metric, than a real cohort differs from unconditional
DDPM output. A reference assembled by the most convenient available route ---
the first $N$ filenames in sorted order --- selects a block that is
systematically atypical of the cohort and inflates the reported score by nearly
a factor of two relative to a subject-diverse reference of the same size. Neither effect is visible in the
number itself; both are invisible in a paper that reports $N$ but not the number
of subjects behind it.

This observation has a constructive consequence. The significance machinery
usually attached to kernel two-sample tests asks whether the real and generated
samples were drawn from the same distribution. For any competent medical
generator that null is false, and in every configuration we tested the test
rejected it at overwhelming significance regardless of generator quality; we measure
permutation $z$-scores above $400$ for every generator we consider, including
one that is a different imaging modality. The comparison a practitioner needs is
different: \emph{is this generator further from the real reference than two real
patient cohorts are from each other?} That question has a natural empirical
null, obtained by measuring distances between disjoint real subject blocks. We
formalise this inter-cohort reference distribution, show that it orders
generators the permutation test cannot, and show that it makes the cohort-size dependence of
any reported score explicit.

With a defensible reference protocol in place, we then address a question that
the shift to domain-pretrained encoders raises but has not settled: a
transformer has many layers, and papers using such encoders as metric backbones
generally read the last one. We audit all twelve blocks of RadioDINO-s16
\citep{perezrad} against four evaluation tasks --- distributional
fidelity, near-duplicate (memorization) detection, mode-drop coverage, and
per-image novelty scoring --- and find that depth is strongly task-dependent and
that the final block is the best choice for none of them. We also find that one
widely used diagnostic, $k$-NN precision/recall
\citep{kynkaanniemi2019pr}, fails its own validation at every depth, in a
direction explained by neighbourhood-radius inflation.

\paragraph{Contributions}
\begin{enumerate}[leftmargin=*, topsep=2pt, itemsep=2pt]
\item We quantify reference-cohort composition as a confound in distributional
      evaluation of medical generative models, and show that on BraTS
      real-versus-real cohort distance exceeds real-versus-generated distance
      (Section~\ref{sec:res_cohort}).
\item We propose an empirical inter-cohort reference distribution that
      complements the permutation test with a comparison against inter-cohort
      heterogeneity,
      and report how it scales with cohort size (Section~\ref{sec:res_null}).
\item We provide a depth-resolved audit of a radiology foundation model across
      four evaluation tasks, showing that the a-priori final-layer choice is
      suboptimal for all of them (Section~\ref{sec:res_depth}).
\item We report that $k$-NN coverage is invalid at every depth in this feature
      space, with a mechanism, and that per-image scoring is more sensitive to
      the choice of distance than to depth (Sections~\ref{sec:res_coverage}
      and~\ref{sec:res_perimage}).
\item We audit ten published medical generation papers and find that none
      reports the number of distinct subjects behind its real reference set,
      which is what makes the confound invisible in practice
      (Section~\ref{sec:res_audit}).
\item We release code implementing the protocol and all experiments.
\end{enumerate}

% =============================================================================
\section{Related work}
\label{sec:related}
% =============================================================================

\paragraph{Distributional metrics and their estimators}
FID models both feature sets as Gaussians and reports the Fr\'echet distance
between them \citep{heusel2017gans}. The estimator is biased at finite sample
size, and the bias depends on $N$ in a way that makes values from different
studies incomparable \citep{chong2020effectively}. KID
\citep{binkowski2018demystifying} replaces the Gaussian assumption with an
unbiased kernel MMD$^2$ estimator \citep{gretton2012kernel}, which removes the
sample-size bias in the estimator but not the dependence on what the reference
set contains. CMMD \citep{jayasumana2024cmmd} keeps the MMD estimator and
substitutes CLIP features for Inception features. Our estimator is a standard
unbiased multi-bandwidth RBF MMD$^2$; the contribution here is not the
estimator but what it is calibrated against.

\paragraph{Domain-specific embeddings}
Several groups have argued that ImageNet features are inappropriate for medical
evaluation and proposed radiology-pretrained substitutes, including
RadImageNet-based FID variants \citep{dohmen2025similarity} and self-supervised
chest and brain encoders \citep{perezrad}. \citet{woodland2024feature}
compare feature extractors for medical FID and report that the choice changes
conclusions. The Fr\'echet Radiomic Distance \citep{konz2025frd} moves in a
different direction, replacing learned features with handcrafted radiomic
descriptors. This literature establishes that embedding choice matters; it does
not examine the reference set, which is what we address.

\paragraph{Beyond a single scalar}
Precision/recall decompositions \citep{sajjadi2018assessing,
kynkaanniemi2019pr} and $\alpha$-precision/$\beta$-recall
\citep{alaa2022faithful} separate sample quality from coverage, and
memorization audits ask whether a generator reproduces training examples
\citep{carlini2023extracting, akbar2025beware}. These diagnostics are typically
computed in the same embedding as the fidelity score and at a single, unexamined
depth. We evaluate them across depth and find that one of them does not
transfer to this feature space at all.

\paragraph{Dependence on the reference set}
The observation that a distributional score depends on what it is measured
against is not new in the natural-image literature, and we do not claim it as
ours. \citet{chong2020effectively} showed that the FID estimator is biased in
the reference sample size, and proposed extrapolating to an infinite-sample
value. Most directly, \citet{lee2026rethinking} study the \emph{geometry} of
the reference dataset, showing across six datasets that distributional density
and effective rank explain much of how FID responds as sample quality improves,
and that concentrated references produce more favourable FID trends than
dispersed ones. Their recommendation --- report distributional metrics together
with descriptors of the reference set --- is one we endorse and extend.

Our question is different in a way that matters for medical imaging, and it is
worth stating precisely which axis is new here. That literature treats the
reference as a set of independent images and asks how its geometry in feature
space affects the metric. We ask what happens when the reference is \emph{not}
a set of independent images: when it is a set of slices drawn from a much
smaller number of patients, so that its effective sample size is governed by the
subject count rather than the image count, and when which subjects are drawn
carries systematic acquisition structure. That distinction has three
consequences with no counterpart in the natural-image setting --- a real-versus-real
comparison that can exceed the real-versus-generated one, a null that shrinks
with subject diversity while the reported score barely moves, and a
subject-level construction protocol --- and it is those we contribute. The
sampling-frame question of whether the reference cohort represents the intended
deployment population is a fourth issue, which neither that literature nor this
paper addresses.

\paragraph{Dataset bias and site effects}
That multi-site medical datasets carry scanner and protocol signatures strong
enough to be recovered by a classifier is well documented
\citep{glocker2019machine, wachinger2021detect}, as is the resulting need for
subject-level rather than image-level data splits
\citep{varoquaux2022machine}. To our knowledge this literature has not been
connected to generative evaluation metrics, where the reference set plays the
role that a test split plays in discriminative work, and where the same
leakage and composition concerns therefore apply.

% =============================================================================
\section{Materials and methods}
\label{sec:methods}
% =============================================================================

\subsection{Data}
\label{sec:data}

\paragraph{Real reference (BraTS)}
We use axial slices extracted from the BraTS 2021 training collection
\citep{baid2021rsna}, giving $38{,}781$ slices from $1{,}251$ subjects
(median $31$ slices per subject) stored as $256\times256$ 8-bit PNG. Slices are
extracted from the FLAIR volume, intensity-normalised per volume to the
$[0,1]$ range after clipping at the $1$st and $99$th percentiles, and retained
where the brain mask covers at least $5\%$ of the frame. Subject identity is
recoverable from the filename, which is what makes the cohort analysis in this
paper possible; we note that a reference distributed as a flat directory of
anonymised slice images would not permit it.

\paragraph{Generated and shifted sets}
Four comparison sets of $500$ images each are used throughout:
\emph{(i)}~\textbf{DDPM}, an unconditional denoising diffusion probabilistic
model \citep{ho2020denoising} with a UNet backbone initialised from
\texttt{google/ddpm-celebahq-256} and fine-tuned on the BraTS slices for $50$
epochs at $256\times256$ (batch $16$, learning rate $10^{-5}$, fp16), sampled
with the standard $1000$-step reverse process;
\emph{(ii)}~\textbf{WDM-3D (MRI)}, axial slices from a 3D wavelet diffusion
model \citep{friedrich2024wdm} trained on BraTS, a same-modality generator of
visibly lower slice-level quality;
\emph{(iii)}~\textbf{LIDC CT}, axial slices from the same 3D architecture
trained on LIDC-IDRI chest CT \citep{armato2011lidc}, a cross-modality shift;
and \emph{(iv)}~\textbf{Retinal}, colour fundus images from RetinaMNIST
\citep{yang2023medmnist}, a far-domain shift. Sets (ii)--(iv) provide an
a-priori severity ordering against which a metric can be checked.

\paragraph{Lesion-level evaluation (LGG)}
For the lesion-specificity analysis we use the lower-grade glioma collection of
\citet{buda2019association}, which supplies expert tumour masks registered to
the slices, allowing perturbation to be confined to annotated tumour tissue or
to area-matched healthy control regions.

\subsection{Reference set construction}
\label{sec:refproto}

Let $\mathcal{S}$ be the set of subjects and $\mathcal{I}_s$ the slices of
subject $s$. A reference of $N$ slices is specified by a subject block
$B \subseteq \mathcal{S}$ and a draw from $\bigcup_{s \in B}\mathcal{I}_s$. We
distinguish three protocols:

\begin{description}[leftmargin=*, topsep=2pt, itemsep=2pt]
\item[\textsc{head}] the first $N$ filenames in sorted order. This is the
      protocol a directory listing yields by default. Because BraTS filenames
      are ordered by subject identifier, it selects the lowest-numbered
      subjects and, at $N=500$, spans only $17$ of them.
\item[\textsc{balanced}$(m)$] $m$ subjects sampled uniformly at random, with
      slices drawn round-robin --- one per subject before any second --- and
      shuffled within each subject. This fixes $N$ while varying subject
      diversity, and avoids drawing only low-index slices, which are
      volume-edge positions and anatomically atypical.
\item[\textsc{disjoint}] two \textsc{balanced} blocks with no shared subject,
      used wherever a real-versus-real comparison is required.
\end{description}

Unless stated otherwise, results use \textsc{balanced}$(100)$ for the reference
and \textsc{disjoint} blocks for real-versus-real comparisons. Subject-level
disjointness matters: a real-versus-real null built by splitting slices
\emph{within} subjects measures near-duplicate distance, not cohort distance,
and is optimistically small.

\subsection{Feature extraction}
\label{sec:features}

We embed images with RadioDINO-s16 \citep{perezrad}, a ViT-S/16
pretrained on radiological images by self-distillation, held frozen throughout.
Images are resized to $224\times224$, replicated to three channels, and
normalised with the encoder's training statistics. For each of the $L=12$
transformer blocks we take the CLS token of that block's output, giving twelve
$384$-dimensional representations per image from a single forward pass. Reading
every depth costs no more than reading one, which is what makes the audit in
Section~\ref{sec:res_depth} inexpensive.

We write $\phi_\ell(x) \in \mathbb{R}^{384}$ for the block-$\ell$ embedding.
Features are $L_2$-normalised before any distance computation.

\subsection{Distributional discrepancy}
\label{sec:mmd}

For real features $X = \{\phi_\ell(x_i)\}_{i=1}^{m}$ and generated features
$Y = \{\phi_\ell(y_j)\}_{j=1}^{n}$ we use the unbiased MMD$^2$ estimator
\citep{gretton2012kernel}
\begin{equation}
\widehat{\mmd}^2(X,Y) =
  \frac{1}{m(m\!-\!1)}\!\sum_{i \neq i'}\! k(x_i,x_{i'})
+ \frac{1}{n(n\!-\!1)}\!\sum_{j \neq j'}\! k(y_j,y_{j'})
- \frac{2}{mn}\sum_{i,j} k(x_i,y_j),
\label{eq:mmd}
\end{equation}
with a multi-bandwidth Gaussian kernel
$k(u,v) = \frac{1}{5}\sum_{c} \exp\!\big(-\|u-v\|^2 / (2 c\,\sigma^2_{\mathrm{med}})\big)$,
$c \in \{\tfrac14, \tfrac12, 1, 2, 4\}$, where $\sigma^2_{\mathrm{med}}$ is the
median pairwise squared distance on the pooled sample. Averaging over
bandwidths avoids committing to a single kernel scale; we report the
single-bandwidth comparison in Section~\ref{sec:res_depth}. The estimator is
unbiased, so values at different $N$ are comparable in expectation --- which is
precisely why the residual dependence on \emph{which} images form the reference
is worth isolating.

\subsection{An empirical inter-cohort reference distribution}
\label{sec:cohortnull}

The conventional null for Eq.~\eqref{eq:mmd} is
$H_0^{\mathrm{perm}}: P_{\mathrm{real}} = P_{\mathrm{gen}}$, tested by pooling
the two samples, relabelling at random, and recomputing. We report it for
comparison, but it is not the hypothesis of interest: no generator satisfies it,
and rejecting it says nothing about whether a generator is usable.

We therefore define a null from the data itself. Partition $\mathcal{S}$ into
$K$ disjoint blocks $B_1, \dots, B_K$ of $m$ subjects each, and draw $N$ slices
from each by the \textsc{balanced} protocol. The
\emph{cohort-heterogeneity null at cohort size $m$} is the set of pairwise
distances
\begin{equation}
\mathcal{N}_m = \big\{\, \widehat{\mmd}^2(R_a, R_b) \;:\; 1 \le a < b \le K \,\big\},
\label{eq:null}
\end{equation}
where $R_a$ is the draw from block $B_a$. For a generated set $G$ scored
against reference $R_1$, we report the cohort-calibrated statistic
\begin{equation}
z_{\mathrm{cohort}} =
  \frac{\widehat{\mmd}^2(R_1, G) - \operatorname{mean}\mathcal{N}_m}
       {\operatorname{sd}\mathcal{N}_m},
\qquad
p_{\mathrm{cohort}} = \frac{1 + |\{v \in \mathcal{N}_m : v \ge \widehat{\mmd}^2(R_1,G)\}|}
                           {1 + |\mathcal{N}_m|}.
\label{eq:cohortz}
\end{equation}
$\mathcal{N}_m$ depends on $m$, so $z_{\mathrm{cohort}}$ is reported together
with the cohort size that produced it. This is a feature rather than a
nuisance: it states in the metric itself how much resolution a reference of that
size affords.

\paragraph{On confidence intervals}
Under $H_0^{\mathrm{perm}}$ the statistic in Eq.~\eqref{eq:mmd} is a degenerate
$U$-statistic, and the ordinary nonparametric bootstrap is not consistent for
it \citep{arcones1992bootstrap}; an $m$-out-of-$n$ bootstrap or subsampling is
required. We therefore never place a bootstrap interval on a null comparison.
Every interval reported here is computed under the alternative --- real against
generated, where the statistic is non-degenerate and the ordinary bootstrap is
consistent --- using $500$ resamples of each sample independently. Null
behaviour is characterised by the empirical distribution $\mathcal{N}_m$
instead, which requires no asymptotic argument.

\subsection{Diagnostics under audit}
\label{sec:diagnostics}

Each diagnostic below is computed at every block $\ell \in \{1,\dots,12\}$ from
the same cached features, and scored by a figure of merit that does not presume
which depth should win.

\paragraph{Distributional fidelity}
Figure of merit: the permutation $z$-score of $\widehat{\mmd}^2(R, G)$ against a
label-shuffling null, which measures how many null standard deviations separate
the observed discrepancy from chance. We additionally check whether each depth
recovers the a-priori severity ordering
$\text{real} < \text{DDPM} < \text{WDM-3D} < \text{CT}$.

\paragraph{Memorization}
A known fraction $r$ of the generated set is replaced by copies of real images
perturbed with Gaussian jitter of scale $\sigma$. The detector flags a generated
image when its nearest-neighbour distance to the reference falls below the
median within-reference nearest-neighbour distance. Figure of merit: the mean
absolute error between the flagged rate and $r$, over
$r \in \{0, 0.1, 0.2, 0.3, 0.4\}$. Sweeping $\sigma$ matters, because at very
small $\sigma$ the task saturates and every depth succeeds.

\paragraph{Coverage}
The generated set is truncated to simulate mode drop, retaining a fraction
$1-d$ for $d \in \{0, 0.2, 0.4, 0.6, 0.8\}$, and $k$-NN precision/recall
\citep{kynkaanniemi2019pr} is computed with $k=5$. Figure of merit: the
Spearman correlation between $d$ and recall, which must be strongly
\emph{negative} for the diagnostic to be usable --- removing diversity must
reduce measured coverage.

\paragraph{Per-image novelty scoring}
Four rules score a single image against the reference: distance to the
reference centroid; $1$-NN and $5$-NN distance; and a Mahalanobis distance using
a covariance estimate shrunk toward a scaled identity
($\hat\Sigma_\varepsilon = (1-\varepsilon)\hat\Sigma + \varepsilon\,\bar\lambda I$,
$\varepsilon = 0.1$), which is necessary because a $384$-dimensional covariance
estimated from $500$ samples is ill-conditioned. Figure of merit: ROC-AUC
separating held-out real images from a positive class. We report two positive
classes of unequal difficulty --- cross-modality CT and the DDPM's own output ---
because a task that saturates cannot distinguish rules or depths.

\paragraph{Degradation monotonicity}
Gaussian noise of increasing $\sigma$ is applied to the generated set. Figure of
merit: Spearman correlation between $\sigma$ and $\widehat{\mmd}^2$, which
should be $+1$.

\subsection{Baselines and implementation}
\label{sec:baselines}

FID and KID are computed with \texttt{torchmetrics} $1.9.0$ at its default
$2048$-dimensional InceptionV3 pool3 features, from \texttt{uint8} inputs at
$256\times256$, letting the library perform its own resize and normalisation.
We flag a pitfall we encountered: a second, hand-rolled FID built on
\texttt{torchvision}'s ImageNet-weight InceptionV3 with ImageNet normalisation
produces values on a different scale for the same images, and mixing the two
within one study yields tables that disagree by tens of FID points while
appearing to describe the same comparison. All FID and KID numbers in this paper
come from one implementation. CMMD \citep{jayasumana2024cmmd} uses CLIP
ViT-L/14 features with the authors' unbiased Gaussian MMD estimator.

Experiments run on a single NVIDIA RTX 4070 Laptop GPU. Extracting all twelve
depths for $500$ images takes $2.5$\,s; one MMD$^2$ evaluation at $N=500$ takes
$4.5$\,s including feature extraction. Building a cohort null from $24$ blocks
requires $24$ extractions and $276$ MMD evaluations. All experiments use seed
$42$ unless a seed sweep is reported.

% =============================================================================
\section{Results}
\label{sec:results}
% =============================================================================

\subsection{Reference cohort composition changes the reported score}
\label{sec:res_cohort}

\paragraph{The default protocol selects an atypical block}
Taking the first $500$ filenames in sorted order from the BraTS slice directory
yields images from $17$ subjects. Randomly sampling $500$ slices from the same
directory spans $413$. Both satisfy a caption reading ``$N = 500$ real
slices'', and they give different answers: MMD$^2$ at \Lb{12} against the same
DDPM samples is $0.150$ for the \textsc{head} reference and $0.077$ for the
random one, a factor of $1.95$.

Subject count alone does not explain this. Table~\ref{tab:blocks} compares
$17$-subject blocks drawn from different positions in the cohort, each
supplying $500$ slices drawn identically. Every block other than the first
gives $\widehat{\mmd}^2 \approx 0.082$; the first gives $0.152$. The effect is
therefore specific to \emph{which} subjects the protocol selects, not to how
many, and it is not a sample-size artefact.

We can state what we measured but not, from this data, its cause. BraTS 2021
does not release a per-subject site label, so we cannot confirm that the
lowest-numbered subjects share a contributing institution; acquisition-site or
protocol grouping correlated with subject-identifier order is the most
economical explanation for a block-structured effect of this size, but it
remains a hypothesis. What matters for the argument here is weaker and
established directly: subject identifiers are not exchangeable, so any protocol
that selects subjects by identifier order is selecting a structured subset of
the cohort.

\begin{table}[t]
\centering
\caption{Seventeen-subject reference blocks drawn from different positions in
the BraTS subject list, each supplying $500$ slices, scored against the same
DDPM sample. Only the block that the default protocol selects is anomalous.}
\label{tab:blocks}
\begin{tabular}{lc}
\toprule
Reference block ($17$ subjects, $N = 500$) & $\widehat{\mmd}^2$ at \Lb{12} \\
\midrule
First $17$ subjects (what \textsc{head} selects) & $0.152$ \\
Middle $17$ subjects                             & $0.084$ \\
Last $17$ subjects                               & $0.082$ \\
Random $17$ subjects (seed $0$)                  & $0.081$ \\
Random $17$ subjects (seed $2$)                  & $0.082$ \\
\bottomrule
\end{tabular}
\end{table}

\paragraph{Real-versus-real distance can exceed real-versus-generated distance}
Table~\ref{tab:realreal} places that site effect on the same scale as the
quantity the metric is meant to report. Two arbitrary $17$-subject real cohorts
differ by $\widehat{\mmd}^2 = 0.017$. The first-$17$ block differs from the rest
of the real cohort by $0.096$. A subject-diverse real reference differs from
DDPM output by $0.080$. The distance between two sets of \emph{real} BraTS
images therefore exceeds the distance between real images and generated ones,
whenever one of them is the atypical block that the default reference protocol
happens to select. We stress the scope of this claim: it is
not that typical real cohort pairs exceed the generated distance --- they do
not, at $0.017$ --- but that the worst-case pair, which is the one the
convenient protocol selects, does.

\begin{table}[t]
\centering
\caption{Distances on a common scale (\Lb{12}, $N = 500$, all cohorts
$17$ subjects so that cohort size is not a confound). The third row involves no
generated images at all, yet exceeds the second.}
\label{tab:realreal}
\begin{tabular}{llc}
\toprule
Comparison & Content & $\widehat{\mmd}^2$ \\
\midrule
Real cohort A vs.\ real cohort B     & real vs.\ real      & $0.017$ \\
Real cohort A vs.\ DDPM              & real vs.\ generated & $0.080$ \\
First-$17$ block vs.\ real cohort A  & real vs.\ real      & $\mathbf{0.096}$ \\
First-$17$ block vs.\ DDPM           & real vs.\ generated & $0.152$ \\
\bottomrule
\end{tabular}
\end{table}

\paragraph{Cohort diversity mostly moves the null, not the signal}
Holding $N = 500$ fixed and sweeping the number of contributing subjects
(Table~\ref{tab:refsweep}, Fig.~\ref{fig:refsweep}) separates two effects. The
real-versus-generated distance is comparatively stable, rising by $9\%$ at
\Lb{12} between a $400$-subject and a $17$-subject reference. The
real-versus-real null is not: it rises from $0.0000$ to $0.0072$ over the same
range. FID moves more than either MMD variant, from $67.9$ to $85.3$
($1.26\times$), consistent with its estimator depending on how well the
reference covariance is estimated. The practical consequence is that reducing
cohort diversity mainly destroys the evaluator's ability to distinguish a real
cohort from a generated one, rather than changing the headline number by much
--- a failure a single reported scalar cannot expose.

\begin{table}[t]
\centering
\caption{Reference cohort diversity at fixed $N = 500$ slices, mean $\pm$ s.d.\
over three draws. Subject diversity moves the real-versus-real null far more
than it moves the real-versus-generated score.}
\label{tab:refsweep}
\begin{tabular}{rccccc}
\toprule
Subjects & $\widehat{\mmd}^2$ \Lb{8} & $\widehat{\mmd}^2$ \Lb{12}
         & null \Lb{8} & null \Lb{12} & FID \\
\midrule
$17$  & $0.1021 \pm 0.0046$ & $0.0801 \pm 0.0030$ & $0.0082$ & $0.0072$ & $85.3 \pm 1.8$ \\
$25$  & $0.1026 \pm 0.0037$ & $0.0799 \pm 0.0025$ & $0.0088$ & $0.0074$ & $82.8 \pm 1.6$ \\
$50$  & $0.0975 \pm 0.0035$ & $0.0751 \pm 0.0026$ & $0.0024$ & $0.0024$ & $71.7 \pm 1.0$ \\
$100$ & $0.0958 \pm 0.0016$ & $0.0739 \pm 0.0005$ & $0.0011$ & $0.0010$ & $72.5 \pm 0.4$ \\
$200$ & $0.0966 \pm 0.0032$ & $0.0749 \pm 0.0023$ & $0.0003$ & $0.0002$ & $68.4 \pm 0.6$ \\
$400$ & $0.0957 \pm 0.0013$ & $0.0736 \pm 0.0008$ & $0.0000$ & $0.0000$ & $67.9 \pm 0.4$ \\
\midrule
$400 \rightarrow 17$ & $\times 1.07$ & $\times 1.09$
  & $+0.0082$ & $+0.0072$ & $\times 1.26$ \\
\bottomrule
\end{tabular}
\end{table}

\begin{figure}[t]
\centering
\includegraphics[width=\linewidth]{Images/reference_set_construction.png}
\caption{Reference cohort diversity at fixed $N = 500$. (a) RadioDINO-s16
MMD$^2$ against DDPM output (solid) and against a held-out real probe
(dashed). The generated distance is comparatively flat, whereas the
real-versus-real null grows sharply as diversity falls. (b) FID on the same
references.}
\label{fig:refsweep}
\end{figure}

\subsection{An empirical inter-cohort reference distribution}
\label{sec:res_null}

Table~\ref{tab:cohortnull} reports the null of Eq.~\eqref{eq:null} at three
cohort sizes, together with each generator's position within it. Two
observations follow.

First, the permutation test does not discriminate here. It assigns $z > 340$ and
$\hat{p}$ at the resolution floor to every generator at every cohort size,
including generators as different from one another as a fine-tuned in-domain
DDPM and colour fundus photography. It answers a question whose answer was
never in doubt.

Second, and more revealing, the two nulls move in \emph{opposite} directions as
the reference improves. Going from a $17$- to a $100$-subject reference, the
inter-cohort statistic for the DDPM rises from $z_{\mathrm{cohort}} = 15.7$
to $89.0$, correctly reporting that a more representative reference resolves
the generator gap more sharply. Over the same range the permutation statistic
\emph{falls}, from $z = 423$ to $z = 347$. The permutation null is built by
shuffling labels within the pooled sample, so a more heterogeneous real set
inflates the null variance and depresses the statistic: a better reference
makes the conventional test report less significance. The inter-cohort reference
distribution does not have this defect, because it is constructed from real
cohorts only.

For an interpretable effect size we report the generator distance as a multiple
of the \emph{mean} real-versus-real cohort distance. With a $17$-subject
reference the DDPM sits at $5.5\times$ the mean real cohort pair; at $50$
subjects, $14.5\times$; at $100$ subjects, $30.3\times$.

We use the mean rather than the maximum deliberately. A ratio to the largest
observed cohort distance is a ratio to an order statistic, whose expectation
grows with the number of pairs compared, so it is not comparable across
conditions that contribute different numbers of pairs. Our $17$- and
$50$-subject conditions both use $K = 24$ blocks and therefore exactly $276$
pairs, and the effect is already established between them with the pair count
held fixed: the ratio to the mean rises from $5.5\times$ to $14.5\times$, and
the ratio to the maximum from $2.3\times$ to $6.9\times$. The $100$-subject
condition admits only $K = 12$ disjoint blocks at this cohort size and so rests
on $66$ pairs; its maximum is drawn from a smaller sample and we do not compare
it to the others on that basis. Table~\ref{tab:cohortnull} reports the mean,
the $95$th percentile and the maximum so the reader can apply whichever
summary they prefer.

\begin{table}[t]
\centering
\caption{Inter-cohort reference distribution at \Lb{12} and generator positions
within it. $z_{\mathrm{cohort}}$ rises with a better reference; the permutation
$z$ falls. ``$\times$mean'' and ``$\times$max'' are the generator distance as a
multiple of the mean and the largest real-versus-real cohort distance; the
former is preferred because the latter is an order statistic whose expectation
depends on the number of pairs. The $17$- and $50$-subject rows rest on the
same $K = 24$ blocks and $276$ pairs and are directly comparable; the
$100$-subject rows rest on $66$ pairs.}
\label{tab:cohortnull}
\small
\begin{tabular}{llcccccc}
\toprule
Cohort & Set & $\widehat{\mmd}^2$ & $\times$mean & $\times$max
 & $z_{\mathrm{cohort}}$ & $z_{\mathrm{perm}}$ & $\hat{p}_{\mathrm{perm}}$ \\
\midrule
\multirow{5}{*}{\shortstack[l]{$17$ subj.\\$K=24$}}
 & null ($276$ pairs) & $0.017 \pm 0.005$ & --- & --- & --- & --- & --- \\
 & DDPM      & $0.095$ & $5.5$  & $2.3$  & $15.7$ & $423.3$ & $<0.002$ \\
 & WDM-3D    & $0.225$ & $13.0$ & $5.5$  & $42.0$ & $709.2$ & $<0.002$ \\
 & LIDC CT   & $0.410$ & $23.6$ & $10.1$ & $79.3$ & $719.1$ & $<0.002$ \\
 & Retinal   & $0.475$ & $27.4$ & $11.7$ & $92.4$ & $810.2$ & $<0.002$ \\
\midrule
\multirow{5}{*}{\shortstack[l]{$50$ subj.\\$K=24$}}
 & null ($276$ pairs) & $0.006 \pm 0.002$ & --- & --- & --- & --- & --- \\
 & DDPM      & $0.081$ & $14.5$ & $6.9$  & $49.5$  & $381.6$ & $<0.002$ \\
 & WDM-3D    & $0.216$ & $38.8$ & $18.5$ & $138.7$ & $689.4$ & $<0.002$ \\
 & LIDC CT   & $0.402$ & $72.2$ & $34.4$ & $261.0$ & $711.1$ & $<0.002$ \\
 & Retinal   & $0.461$ & $82.7$ & $39.4$ & $299.4$ & $810.4$ & $<0.002$ \\
\midrule
\multirow{5}{*}{\shortstack[l]{$100$ subj.\\$K=12$}}
 & null ($66$ pairs) & $0.003 \pm 0.001$ & --- & --- & --- & --- & --- \\
 & DDPM      & $0.079$ & $30.3$  & $15.9$ & $89.0$  & $347.2$ & $<0.002$ \\
 & WDM-3D    & $0.220$ & $84.3$  & $44.1$ & $253.3$ & $631.4$ & $<0.002$ \\
 & LIDC CT   & $0.403$ & $154.4$ & $80.8$ & $466.0$ & $678.3$ & $<0.002$ \\
 & Retinal   & $0.463$ & $177.2$ & $92.7$ & $535.6$ & $766.2$ & $<0.002$ \\
\bottomrule
\end{tabular}
\end{table}

\begin{figure}[t]
\centering
\includegraphics[width=\linewidth]{Images/cohort_heterogeneity_null.png}
\caption{Generator distances against the real cohort-heterogeneity null at
three cohort sizes. Histograms are pairwise distances between disjoint real
patient blocks; vertical lines mark generators. The horizontal axis is
logarithmic, because the null lies roughly two orders of magnitude below the
generator distances and is not visible on a linear scale. As cohorts grow the
null shifts left and concentrates, while the generator distances stay
essentially fixed --- the separation the evaluation can resolve is set by the
reference, not by the generators.}
\label{fig:cohortnull}
\end{figure}

\subsection{What do published papers actually report?}
\label{sec:res_audit}

Sections~\ref{sec:res_cohort}--\ref{sec:res_null} show that reference
composition changes the reported number. A fair objection is that we
constructed a deliberately poor selection rule and then demonstrated that it is
poor. Whether the finding matters depends on an empirical question about the
field: do published papers report enough about their reference set for a reader
to know which rule was used?

We audited ten papers that evaluate medical image generation, spanning brain
MRI, chest X-ray, cardiac MRI and CT, from 2022 to 2025, and including two
papers whose subject is evaluation methodology rather than generation. For each
we recorded five properties answerable from the text: whether the size of the
real reference set is stated; whether the number of distinct subjects behind it
is stated; whether the selection procedure is described; whether train/test
splits are stated to be subject-level; and whether the reference is drawn from
the generator's training data. Table~\ref{tab:audit} gives the result.

\begin{table}[t]
\centering
\caption{Reference-set reporting in ten medical image generation papers.
\cmark~= stated, $\sim$~= partially stated, \xmark~= not stated,
``?''~= cannot be determined from the text. One paper uses no distributional
metric and is marked --- in the metric-specific columns. Percentages exclude
that paper.}
\label{tab:audit}
\footnotesize
\begin{tabular}{llccccc}
\toprule
Paper & Metric & $N$ imgs & $N$ subj & Selection & Subj.\ split & Ref.\ from train \\
\midrule
Pinaya et al.\ 2022        & FID       & \xmark & \xmark & \xmark & \xmark & ? \\
Khader et al.\ 2023        & MS-SSIM   & ---    & ---    & ---    & \xmark & --- \\
Akbar et al.\ 2023         & FID       & \xmark & $\sim$ & \xmark & \cmark & ? \\
Kong et al.\ 2023          & FID       & \xmark & \xmark & \xmark & \xmark & ? \\
Woodland et al.\ 2024      & Fr\'echet & $\sim$ & \xmark & $\sim$ & \xmark & ? \\
Hashmi et al.\ 2024        & FID       & \xmark & \xmark & \xmark & \xmark & $\sim$ \\
BrainMRDiff 2025           & MMD       & \xmark & \xmark & \xmark & \cmark & ? \\
FADM 2024                  & FID       & \xmark & $\sim$ & \xmark & \cmark & ? \\
Rajesh et al.\ 2025        & FID, ECS  & \xmark & \xmark & \xmark & \xmark & \cmark \\
Perceptual Eval.\ 2025     & reader    & $\sim$ & \xmark & \xmark & \xmark & ? \\
\midrule
\multicolumn{2}{l}{Fully stated}       & $0/9$  & $0/9$  & $0/9$  & $3/10$ & $1/9$ \\
\bottomrule
\end{tabular}
\end{table}

Three results stand out. \emph{No paper in the sample states how many distinct
subjects its real reference set is drawn from.} No paper fully states the
reference set size, and none describes the selection procedure well enough to
reproduce it. Only three of ten state that their splits are subject-level, and
only one states explicitly whether the reference overlaps the generator's
training data --- and that one does so because the overlap is a deliberate part
of its design.

The consequence is not that these papers are wrong. It is that a reader cannot
distinguish a reference drawn from $17$ subjects from one drawn from $400$,
and Section~\ref{sec:res_cohort} shows those are different instruments giving
answers that differ by a factor of two. The information needed to tell them
apart is, in this sample, never present. That is what makes reference
composition a reporting problem rather than a modelling one, and it is why our
recommendations in Section~\ref{sec:protocol} are about what to report rather
than about what to compute.

\paragraph{Scope of the audit}
This is a purposive sample of ten papers assembled from literature searches on
medical image synthesis with distributional metrics, not a systematic review
with a pre-registered protocol, and it is small enough that the percentages
should be read as indicative rather than as prevalence estimates. Judgements
are made from the main text; a detail buried in a supplement or a code release
would be scored \xmark\ here. We report the per-paper evidence for every
judgement in the released audit record so that individual scores can be
checked and disputed.

\subsection{Encoder depth is task-dependent}
\label{sec:res_depth}

All results in this section use the \textsc{disjoint} reference protocol of
Section~\ref{sec:refproto} at $N = 500$, so the cohort effects of
Section~\ref{sec:res_cohort} are controlled. Table~\ref{tab:depthgrid} and
Fig.~\ref{fig:depthgrid} summarise every task at every block.

\paragraph{Fidelity peaks before the final block}
Permutation $z$ for the real-versus-DDPM comparison rises steeply through the
early blocks, plateaus over \Lb{7}--\Lb{9}, and then declines. The maximum is at
\Lb{8} ($z = 439.3$), with \Lb{7} ($436.9$) and \Lb{9} ($436.6$) statistically
indistinguishable from it. The final block gives $z = 342.4$, which is $22.1\%$
below the plateau. Reading the last layer, the default in work that uses a
foundation model as a metric backbone, therefore discards roughly a fifth of the
available discriminability on this comparison.

The raw discrepancies tell the same story: $\widehat{\mmd}^2$ against the DDPM
is $0.1042$ at \Lb{7} and $0.0732$ at \Lb{12}, while the real-versus-real null
is $\le 0.0003$ at every block. The decline into the final blocks is a loss of
signal, not a change of scale.

Severity ordering, by contrast, is depth-robust. Every one of the twelve blocks
recovers the a-priori ordering
$\text{real} < \text{DDPM} < \text{WDM-3D} < \text{CT} < \text{retinal}$ with
Spearman $\rho = +1.000$. If the only requirement is to rank generators of
clearly different quality, depth does not matter; if the requirement is to
resolve two similar generators, it does.

\paragraph{Near-duplicate detection improves with depth}
Memorization recovery error falls essentially monotonically with depth
(Table~\ref{tab:depthgrid}, second column): mean absolute error is $0.512$ at
\Lb{1}, $0.027$ at \Lb{6}, $0.007$ at \Lb{9} and $0.004$ at \Lb{10}--\Lb{11}.
Shallow blocks are not merely worse, they are unusable: at \Lb{1} the detector
recovers essentially none of the injected duplicates. The behaviour is stable
across jitter scales up to $\sigma = 0.02$ and collapses for all depths at
$\sigma \ge 0.1$, where the perturbation has destroyed the near-duplicate
relationship the detector is meant to find.

This result contradicts an earlier analysis of ours conducted on a
\textsc{head} reference, which reported that shallow blocks outperformed
mid-depth ones for this task. That conclusion was an artefact of the
$17$-subject reference: when the reference spans few subjects, its
within-reference nearest-neighbour distances are small, the detection threshold
is correspondingly tight, and the ranking across depths inverts. We report this
explicitly because it is a concrete instance of the paper's general point ---
the reference protocol changed the sign of a methodological conclusion, not
merely the value of a number.

\paragraph{Degradation monotonicity fails only at the shallowest blocks}
Under a Gaussian noise ladder, $\widehat{\mmd}^2$ increases monotonically with
$\sigma$ at every block from \Lb{3} onwards ($\rho = +1.000$). \Lb{1}
($\rho = +0.857$) and \Lb{2} ($\rho = +0.893$) are non-monotone: mild noise
initially \emph{decreases} the measured discrepancy, because at those depths the
representation is dominated by low-level texture statistics that noise makes
more similar between the two sets before it makes them different. A metric read
at \Lb{1} would report that a slightly corrupted generator had improved.

\begin{table}[t]
\centering
\caption{Depth $\times$ task grid for RadioDINO-s16, $N = 500$, patient-disjoint
reference. Best value per column in bold; the a-priori choice for each task is
marked $\dagger$. No single depth wins every task, and the final block wins
none.}
\label{tab:depthgrid}
\small
\begin{tabular}{rccccc}
\toprule
& Fidelity & Memorization & Coverage & \multicolumn{2}{c}{Per-image AUC} \\
\cmidrule(lr){5-6}
Block & $z_{\mathrm{perm}}$ $\uparrow$ & MAE $\downarrow$
      & $\rho(d,\text{recall})$ $\downarrow$ & centroid & Mahalanobis \\
\midrule
\Lb{1}  & $171.6$ & $0.512$ & $+0.900$ & $0.592$ & $0.461$ \\
\Lb{2}  & $212.6$ & $0.163$ & $+1.000$ & $0.664$ & $0.725$ \\
\Lb{3}  & $185.0$ & $0.180$ & $+1.000$ & $0.646$ & $0.622$ \\
\Lb{4}  & $221.1$ & $0.114$ & $+0.900^{\dagger}$ & $0.654$ & $0.771$ \\
\Lb{5}  & $267.4$ & $0.083$ & $+1.000$ & $0.628$ & $0.858$ \\
\Lb{6}  & $326.5$ & $0.027$ & $+0.900$ & $0.674$ & $0.921$ \\
\Lb{7}  & $436.9$ & $0.013$ & $+1.000$ & $0.744$ & $0.928$ \\
\Lb{8}  & $\mathbf{439.3}$ & $0.016$ & $+0.900$ & $0.796$ & $0.949$ \\
\Lb{9}  & $436.6$ & $0.007^{\dagger}$ & $+0.975$ & $\mathbf{0.846}$ & $\mathbf{0.972}$ \\
\Lb{10} & $404.6$ & $\mathbf{0.004}$ & $+1.000$ & $0.840$ & $\mathbf{0.972}$ \\
\Lb{11} & $368.6$ & $\mathbf{0.004}$ & $+1.000$ & $0.821$ & $0.967$ \\
\Lb{12} & $342.4^{\dagger}$ & $0.008$ & $+0.821$ & $0.810$ & $0.962$ \\
\midrule
Best block & \Lb{8} & \Lb{10}--\Lb{11} & none & \Lb{9} & \Lb{9}--\Lb{10} \\
\bottomrule
\end{tabular}
\end{table}

\begin{figure}[t]
\centering
\includegraphics[width=\linewidth]{Images/depth_task_grid.png}
\caption{Depth profiles for four evaluation tasks. (a) Fidelity
discriminability peaks at \Lb{7}--\Lb{9}, not at the final block. (b)
Memorization recovery error falls with depth; curves are jitter scales, and the
task collapses for all depths at large $\sigma$. (c) $k$-NN coverage: the
Spearman correlation between mode drop and recall is positive at every depth,
the opposite of the required sign. (d) Per-image scoring: the choice of
distance separates the curves more than depth does. Dashed vertical lines mark
the a-priori depth for each task.}
\label{fig:depthgrid}
\end{figure}

\subsection{$k$-NN coverage is invalid at every depth}
\label{sec:res_coverage}

The mode-drop validation asks a minimal question: if diversity is deliberately
removed from the generated set, does measured coverage fall? For $k$-NN recall
\citep{kynkaanniemi2019pr} in this feature space, the answer is no at
every one of the twelve blocks. The Spearman correlation between drop fraction
and recall is \emph{positive} everywhere, ranging from $+0.821$ at \Lb{12} to
$+1.000$ at five blocks (Table~\ref{tab:depthgrid}, third column). Discarding
$80\%$ of the generated samples \emph{increases} measured recall --- at \Lb{9},
from $0.008$ to $0.114$.

The mechanism is visible in the raw values and is a known property of the
estimator rather than a pathology of the encoder. Recall is the fraction of real
points falling inside the union of $k$-NN balls around generated points, and
each ball's radius is that point's own $k$-th nearest-neighbour distance
\emph{within the generated set}. Removing generated samples increases every
surviving point's $k$-NN radius, inflating each ball. Below some density the
inflation of the balls outpaces the loss of centres, and measured coverage
rises. In a $384$-dimensional self-supervised embedding, where the generated set
occupies a thin region and absolute recall is already very low ($0.008$--$0.044$
at $d = 0$ for blocks \Lb{7}--\Lb{12}), the estimator is in that regime from the
start.

We therefore report no coverage number in this feature space and recommend that
others do not either without first running this validation. We note that the
diagnostic may behave correctly in embeddings where the generated set is denser
relative to the real manifold; the claim here is specific to the setting, not a
general refutation of the estimator.

\subsection{Per-image scoring depends more on the distance than on the depth}
\label{sec:res_perimage}

Two per-image tasks were run: separating held-out real slices from
cross-modality CT, and separating them from the DDPM's own output. The first is
saturated --- every rule at every depth achieves AUC $\ge 0.997$ --- and is
therefore useless for comparing rules or depths. We note this because
cross-modality OOD detection is a common headline experiment in this
literature, and a saturated benchmark cannot support the conclusions usually
drawn from it.

On the discriminative same-modality task, the ranking of rules is consistent
across depth and larger than the effect of depth itself
(Table~\ref{tab:depthgrid}, last two columns). At the best depth, shrinkage
Mahalanobis reaches AUC $0.972$, $1$-NN $0.942$, $5$-NN $0.929$, and distance to
the reference centroid $0.846$. The centroid rule is worst at every block from
\Lb{2} onward, by a margin of $9$--$23$ AUC points. This is expected on
geometric grounds: a centroid distance assumes an isotropic, unimodal reference,
whereas a medical cohort is neither, and a shrinkage Mahalanobis distance
accounts for the anisotropy that the raw Euclidean distance ignores. Where a
per-image score is required --- for flagging individual synthetic images before
clinical use, for instance --- the distance should be chosen deliberately rather
than defaulted to a centroid.

\subsection{Cross-metric comparison on a single protocol}
\label{sec:res_baselines}

Table~\ref{tab:unified} reports every metric on one reference protocol and one
implementation per metric, which Section~\ref{sec:res_cohort} argues is
necessary for the numbers to be comparable at all. Two rows contain no
generated images. \textsc{real probe} is a second, subject-disjoint block drawn
by the same \textsc{balanced} procedure as the reference: the value a metric
should return when nothing is wrong. \textsc{real head} is the block that the
default filename-ordered protocol selects --- also entirely real BraTS slices,
at the same resolution, from the same collection, preprocessed identically.

\paragraph{A real reference can look like a generative model}
The two real rows are the practical summary of this paper. Drawn correctly, a
real block sits at $\widehat{\mmd}^2 = 0.002$ and FID $= 36.0$. Drawn by the
convenient protocol, the same kind of comparison --- real against real ---
reaches $\widehat{\mmd}^2 = 0.067$ and FID $= 64.6$. The DDPM, an actual
generative model, sits at $0.079$ and $70.8$. On FID the gap between a misdrawn
real reference and a real generator is about six points, set against the
$29$-point spread between the two real rows. A reader given only the scalar
$64.6$ has no way to tell that no generative model was involved in producing
it. This holds for every metric in the table, including the two
natural-image baselines, so it is a property of the protocol rather than of the
medical backbone.

\paragraph{Ordering, and where the backbones disagree}
All four metrics order the three generated sets in the same way relative to the
reference, and all place the weak same-modality generator closer than the
cross-modality shift (Table~\ref{tab:unified}). The metrics disagree on the
\emph{far} end of the ordering. The RadioDINO MMD places retinal fundus further
from brain MRI ($0.463$ at \Lb{12}) than chest CT ($0.403$); FID and KID place
CT further ($322.2$) than retinal ($248.4$). Neither ordering is wrong in the
abstract --- they encode different notions of similarity. RadioDINO, trained on
radiology, treats CT and MRI as neighbouring modalities and ophthalmic
photography as remote; InceptionV3, trained on natural images, responds to
surface appearance, and a colour fundus photograph is closer to its training
distribution than a greyscale CT slice. For evaluating a CT or MRI generator we
regard the radiological ordering as the more useful of the two, but this is an
argument about which similarity a practitioner wants, not evidence that one
backbone is more accurate.

We had previously reported, on a \textsc{head} reference, that CMMD inverts the
ordering of a weak same-modality generator against a cross-modality CT shift.
That measurement was made on the reference protocol this paper argues against,
and we therefore do not carry the claim forward; whether it survives a
subject-balanced reference is left open.

\begin{table}[t]
\centering
\caption{All metrics on one reference protocol (\textsc{balanced}$(100)$,
$N = 500$) and one implementation each. The first two rows contain only real
images. Bootstrap intervals ($60$ resamples) are given for comparisons under
the alternative; none is given for the real rows, where the statistic
approaches degeneracy and the ordinary bootstrap is not consistent
(Section~\ref{sec:cohortnull}).}
\label{tab:unified}
\small
\begin{tabular}{llcccc}
\toprule
Comparison set & Content & $\widehat{\mmd}^2$ \Lb{8} & $\widehat{\mmd}^2$ \Lb{12}
 & FID & KID \\
\midrule
\textsc{real probe}  & real vs.\ real & $0.0016$ & $0.0015$ & $36.0$  & $0.0009$ \\
\textsc{real head}   & real vs.\ real & $0.0628$ & $0.0674$ & $64.6$  & $0.0299$ \\
\midrule
DDPM           & real vs.\ gen. & $0.1028$ & $0.0792$ & $70.8$  & $0.0380$ \\
WDM-3D (MRI)   & real vs.\ gen. & $0.2155$ & $0.2204$ & $174.8$ & $0.1574$ \\
LIDC CT        & real vs.\ gen. & $0.3886$ & $0.4034$ & $322.2$ & $0.3725$ \\
Retinal        & real vs.\ gen. & $0.5026$ & $0.4632$ & $248.4$ & $0.2851$ \\
\bottomrule
\end{tabular}
\end{table}

\paragraph{A note on the degenerate row}
The bootstrap interval our procedure returns for \textsc{real probe} at \Lb{8}
is $[0.0025, 0.0047]$ for a point estimate of $0.0016$ --- an interval that does
not contain the estimate it is meant to bracket. This is not an implementation
error but the expected behaviour of the ordinary bootstrap on a degenerate
$U$-statistic \citep{arcones1992bootstrap}, and it is precisely why we
characterise null behaviour with the empirical inter-cohort distribution of
Section~\ref{sec:cohortnull} rather than with bootstrap intervals. We report the
pathology rather than suppressing it, because a reader who applies the
bootstrap uncritically to a near-null comparison will meet it.

\subsection{Lesion-level specificity}
\label{sec:res_lesion}

A metric intended for radiology should respond more strongly to changes in
pathological tissue than to equivalent changes elsewhere. We test this on the
LGG collection \citep{buda2019association}, which provides expert tumour masks.
Gaussian noise of scale $\sigma$ is applied either inside the annotated tumour
region or inside an area-matched control region, and we report the ratio of the
induced discrepancy in the two conditions. Two controls are used: a mirrored
region at the contralateral position, and a texture-matched region selected to
have comparable gradient energy ($4.97\times10^{4}$ for tumour,
$4.77\times10^{4}$ mirrored, $5.33\times10^{4}$ texture-matched), which
controls for the possibility that the metric merely responds to local
structural complexity.

The RadioDINO MMD is lesion-selective where an Inception-based MMD is not, but
only in the upper part of the sweep (Table~\ref{tab:lesion}). At $\sigma = 60$
and $\sigma = 90$ the bootstrap interval excludes unity ($1.57$,
$[1.03, 1.97]$ and $1.39$, $[1.06, 1.71]$ against the mirrored control); at
$\sigma = 45$ the point estimate is larger ($1.71$) but the interval
$[0.94, 2.14]$ includes unity. At $\sigma \le 30$ the control-region response is
near zero and the ratio is numerically degenerate, so we report no value. The
Inception baseline stays within $1.01$--$1.08$ throughout, with intervals
covering unity at every $\sigma$.

We state this as suggestive rather than established, for two reasons. The
effect is individually significant at only two of the six perturbation levels
tested. More importantly, the perturbation is additive Gaussian noise confined
to a mask, which is a proxy for pathology-relevant change rather than an
instance of it: it establishes that the embedding weights tumour regions
differently from matched healthy regions, but not that it tracks clinically
meaningful pathology. Establishing the latter would require
pathology-preserving edits, lesion insertion or removal, or agreement with
radiologist judgement, none of which we attempt here.

\begin{table}[t]
\centering
\caption{Lesion selectivity: ratio of the discrepancy induced by perturbing
tumour tissue to that induced by perturbing an area-matched control region.
A value of $1.0$ means no lesion selectivity. Bootstrap $95\%$ intervals;
$\dagger$ marks intervals excluding unity. Values at $\sigma \le 30$ are
omitted because the control response is near zero and the ratio is degenerate.}
\label{tab:lesion}
\small
\begin{tabular}{rcccc}
\toprule
& \multicolumn{2}{c}{RadioDINO MMD} & \multicolumn{2}{c}{Inception MMD} \\
\cmidrule(lr){2-3}\cmidrule(lr){4-5}
$\sigma$ & mirrored & texture-matched & mirrored & texture-matched \\
\midrule
$45$ & $1.71\,[0.94, 2.14]$ & $1.81\,[0.94, 2.12]$
     & $1.08\,[0.92, 1.24]$ & $1.03\,[0.87, 1.16]$ \\
$60$ & $1.57\,[1.03, 1.97]^{\dagger}$ & $1.59\,[1.02, 1.90]^{\dagger}$
     & $1.08\,[0.93, 1.22]$ & $1.05\,[0.91, 1.17]$ \\
$90$ & $1.39\,[1.06, 1.71]^{\dagger}$ & $1.40\,[1.03, 1.69]^{\dagger}$
     & $1.06\,[0.93, 1.19]$ & $1.06\,[0.94, 1.17]$ \\
\bottomrule
\end{tabular}
\end{table}

\begin{figure}[t]
\centering
\includegraphics[width=0.9\linewidth]{Images/lgg_noise_sweep.png}
\caption{Lesion selectivity across the perturbation sweep on LGG, with
area-matched mirrored and texture-matched controls. Shaded bands are bootstrap
$95\%$ intervals. The shaded region at low $\sigma$ marks where the
control-region response is near zero and the ratio is not measurable.}
\label{fig:lesion}
\end{figure}

% =============================================================================
\section{Discussion}
\label{sec:discussion}
% =============================================================================

\subsection{What a reported metric value does and does not mean}

The results in Section~\ref{sec:res_cohort} put a bound on how much a single
reported number can carry. Two studies can both report ``MMD$^2 = 0.15$ against
$500$ real BraTS slices'' while differing by a factor of two in the quantity
they intend to describe, purely through reference composition. Nothing in the
number, the sample size, or the estimator's unbiasedness reveals the
difference. Unbiasedness is a property of the estimator with respect to a fixed
pair of distributions; it says nothing about whether the empirical reference
represents the population it stands for.

This is the generative-evaluation counterpart of a problem that discriminative
medical imaging settled some time ago. Subject-level rather than image-level
splits are standard there, because slices from one subject leak between train
and test \citep{varoquaux2022machine}, and multi-site datasets are known to
carry recoverable scanner signatures \citep{glocker2019machine}. Generative
evaluation inherited neither convention, because it inherited its protocols from
natural-image benchmarks where neither problem exists. The reference set in a
generative study plays the structural role that a test split plays in a
discriminative one, and it deserves the same care.

\subsection{Why the conventional significance test is the wrong instrument}

Attaching a permutation test to a kernel two-sample statistic is technically
correct and practically uninformative. Its null --- that real and generated
samples come from the same distribution --- is false by construction for every
generative model anyone would submit for evaluation. Our measurements bear this
out: the test assigns $z > 340$ to an in-domain DDPM and to colour fundus
photography alike (Table~\ref{tab:cohortnull}), separating neither from the
other nor from anything else.

The behaviour under improving references is worse than merely uninformative.
Because the permutation null is built by shuffling labels within the pooled
sample, a more heterogeneous real set inflates the null variance and
\emph{depresses} the reported statistic. We say heterogeneous rather than
representative deliberately: we have no sampling frame for the target
population, so a wider subject draw is demonstrably more diverse but not
thereby shown to be more representative. A
practitioner who improves their reference from $17$ to $100$ subjects would see
the conventional significance of their result fall, from $z = 423$ to $z = 347$,
while the evaluation actually became better resolved. The inter-cohort
statistic moves the other way, from $15.7$ to $89.0$.

We do not claim the inter-cohort reference distribution is the only correct construction. It requires
subject labels, which not every released dataset provides; it is defined at a
particular cohort size, so two studies must match cohort size to be comparable;
and with $K$ blocks it rests on $K(K-1)/2$ non-independent pairs, so its tail
should be read as descriptive rather than as a calibrated $p$-value. We report
$z_{\mathrm{cohort}}$ and the ratio to the \emph{mean} observed cohort distance
for that reason, rather than leaning on $p_{\mathrm{cohort}}$, whose resolution
is bounded by the number of pairs, or on a ratio to the maximum, which is an
order statistic that is not comparable across conditions contributing different
numbers of pairs. What we do claim is that a null built from real
cohorts answers the question a practitioner is actually asking, and that the
permutation null does not.

\subsection{Reading depth deliberately}

The depth audit yields a practical result and a cautionary one. The practical
result is that a single frozen encoder supports several evaluation tasks at
different depths, at no additional cost: all twelve representations come from
one forward pass, so reading \Lb{8} for fidelity and \Lb{10} for near-duplicate
detection is free. The cautionary result is that the final block, which is what
the literature reads by default, is the best choice for none of the four tasks
we measured, and is $22\%$ below the fidelity plateau. We name \Lb{8} as the
maximum but not as ``the'' optimal depth: \Lb{7}, \Lb{8} and \Lb{9} are not
distinguishable from one another, and the defensible statement is that the
optimum lies in a mid-to-late plateau, not at the final block.

We are explicit that our depth recommendations are empirical rather than
principled. We do not have a theory that predicts why fidelity discriminability
should peak at \Lb{8} in this encoder, and we would not expect the specific
number to transfer to a different backbone or a different modality. What we
expect to transfer is the methodology: the depth profile is cheap to measure,
the figures of merit above do not presume an answer, and any study using a
foundation model as a metric backbone can run the audit for its own setting in
minutes.

The memorization reversal described in Section~\ref{sec:res_depth} is worth
restating as a general warning. On a \textsc{head} reference we previously
concluded that shallow blocks were better than mid-depth blocks for
near-duplicate detection. On a subject-disjoint reference the ordering reverses
and becomes monotone in the opposite direction. The reference protocol did not
shift a number by a few percent; it inverted a methodological conclusion. Any
depth study, ablation, or backbone comparison conducted on a convenience
reference is exposed to the same risk.

\subsection{Limitations}

\paragraph{One dataset for the cohort analysis}
The cohort-heterogeneity results are measured on BraTS, whose multi-site
composition and subject-ordered filenames make the effect easy to isolate. We
expect the phenomenon wherever a reference is assembled from a multi-site
collection of volumetric studies, but we have not measured its magnitude
elsewhere, and the specific inflation factors reported here should not be
assumed to transfer. Whether datasets curated from a single site show a
correspondingly smaller effect is an open question that the protocol in
Section~\ref{sec:cohortnull} is designed to answer.

\paragraph{One encoder for the depth analysis}
The depth grid uses RadioDINO-s16. The finding that depth is task-dependent is
likely general; the specific optima are not claimed to be.

\paragraph{Generators are not state of the art}
The DDPM and WDM-3D samples serve as objects of measurement, not as
contributions. A stronger generator would sit closer to the cohort null and
would make the resolution argument of Section~\ref{sec:res_null} more pressing,
not less: the closer a generator comes to real data, the more the reference
protocol determines whether the remaining gap can be detected at all.

\paragraph{Lesion-level specificity is a partial result}
Using expert glioma masks with area-matched mirror and texture-matched
controls, the RadioDINO MMD responds $1.4$--$1.8\times$ more strongly to
perturbation confined to tumour tissue than to matched healthy tissue, where an
Inception-based MMD responds at $1.05$--$1.08\times$. The effect is measurable
only in the upper part of the noise sweep: bootstrap intervals exclude unity at
$\sigma = 60$ ($1.57$, $[1.03, 1.97]$) and $\sigma = 90$ ($1.39$,
$[1.06, 1.71]$) but not at $\sigma = 45$ ($1.71$, $[0.94, 2.14]$), and at
$\sigma \le 30$ the ratio is numerically degenerate because the control-region
response is near zero. We report it as suggestive support for the domain
backbone rather than as an established property.

\paragraph{Coverage}
Our coverage result is a negative one about a specific estimator in a specific
embedding. It does not establish that no coverage diagnostic can work here, and
we did not evaluate alternatives such as $\alpha$-precision/$\beta$-recall
\citep{alaa2022faithful} at every depth. The recommendation is to run the
mode-drop validation before reporting any coverage number, not to abandon
coverage as a concept.

% =============================================================================
\section{Recommended protocol}
\label{sec:protocol}
% =============================================================================

The following is what we would ask of a paper reporting a distributional metric
for medical image generation, and what our released code implements by default.

\begin{enumerate}[leftmargin=*, topsep=2pt, itemsep=3pt]
\item \textbf{Report the number of distinct subjects behind every image set,
      not only the number of images.} A reference of $500$ slices from $17$
      subjects and one from $400$ are different instruments. This costs one
      number in a table and makes the remaining recommendations checkable.
\item \textbf{Build references by subject, not by filename.} Sample subjects
      first, then draw slices balanced across them and shuffled within them.
      Sorted-order convenience selects a site.
\item \textbf{Keep real-versus-real comparisons subject-disjoint.} A null built
      by splitting slices within subjects measures near-duplicate distance and
      is optimistically small.
\item \textbf{Contextualise against inter-cohort variation.} Report
      $z_{\mathrm{cohort}}$ and the ratio to the \emph{mean} real-versus-real
      cohort distance, at a stated cohort size, alongside the raw score. Avoid
      ratios to the maximum: it is an order statistic whose expectation grows
      with the number of cohort pairs compared. Report the permutation
      $p$-value if desired, but do not read it as evidence of generator
      quality.
\item \textbf{Audit read-out depth; do not copy ours.} Run the depth audit
      for your own encoder --- it costs one forward pass. Our optimum
      (\Lb{7}--\Lb{9} for fidelity, with those three not distinguishable from
      one another) is a property of this encoder and this data, and should not
      be transplanted.
\item \textbf{Validate every diagnostic before reporting it.} Mode drop must
      reduce measured coverage; injected duplicates must be recovered at the
      injected rate. A diagnostic that fails its own validation should not
      appear in a results table.
\item \textbf{Use one implementation per baseline metric.} FID computed from
      two different Inception variants differs by tens of points on identical
      images.
\end{enumerate}

% =============================================================================
\section{Conclusion}
\label{sec:conclusion}
% =============================================================================

We set out to audit a domain-pretrained embedding for generative medical image
evaluation and found that the more consequential variable lies upstream of the
embedding entirely. On BraTS, how the real reference cohort is assembled changes
the reported discrepancy by up to a factor of two, moves the real-versus-real
null by more than it moves the real-versus-generated score, and --- in the
worst case that the most convenient protocol happens to select --- produces a
distance between two sets of real images that exceeds the distance between real
and generated ones. Reference composition is not a second-order detail in this
setting; it is comparable in magnitude to the effect being measured.

We proposed calibrating against a null built from distances between disjoint
real patient cohorts, which asks whether a generator differs from the reference
by more than routine inter-cohort variation. Unlike the conventional
permutation test, which assigns overwhelming significance to every generator and
becomes \emph{less} significant as the reference improves, the cohort-calibrated
statistic tracks reference quality in the right direction and yields an
interpretable effect size.

Within that protocol, a depth-resolved audit of RadioDINO-s16 shows that
encoder depth is task-dependent and that the final block --- the default
read-out --- is optimal for none of the four tasks measured, falling $22\%$
below the fidelity plateau at \Lb{7}--\Lb{9}, whose members are not
distinguishable from one another. One widely used diagnostic, $k$-NN
precision/recall coverage, fails its own mode-drop validation at every depth
through neighbourhood-radius inflation, and we recommend against reporting it
in this feature space without validation. For per-image scoring, the choice of
distance matters more than depth.

None of these findings requires a new metric. They ask for a more careful
account of what an existing one is measured against.

% =============================================================================
\section*{Code and data availability}
% =============================================================================

Code implementing the metric, the inter-cohort reference distribution, the
depth audit, the literature audit record and
every experiment reported here is available at
\url{https://github.com/Nishan-Charlie/m3-cohort-audit}. BraTS 2021
\citep{baid2021rsna} and the LGG collection \citep{buda2019association} are
available from their respective repositories under their own terms; we
redistribute no patient data. Scripts to regenerate every table and figure from
raw data are included, together with the exact subject lists defining each
cohort block.

\section*{CRediT authorship contribution statement}

\textbf{Sathiyamohan Nishankar:} Conceptualization, Methodology, Software,
Validation, Formal analysis, Investigation, Data curation, Writing --- original
draft, Writing --- review \& editing, Visualization.

\section*{Declaration of competing interest}

The author declares no known competing financial interests or personal
relationships that could have appeared to influence the work reported in this
paper.


\bibliographystyle{elsarticle-harv}
\bibliography{refs}

\end{document}

